\documentclass{article}
\usepackage[T1]{fontenc}
\usepackage{lmodern}
\usepackage{graphicx}
\usepackage{amsmath,amssymb}
\usepackage[a4paper,margin=2cm]{geometry}
\usepackage[absolute,overlay]{textpos}
\setlength{\TPHorizModule}{1mm}
\setlength{\TPVertModule}{1mm}
\usepackage[indonesian]{babel}
\usepackage{hyperref}
\setlength{\emergencystretch}{2em}
% unit: Halaman 46

\begin{document}

% === Halaman 46 (llm) ===

\section*{Perbandingan Elgamal dengan Elliptic Curve Elgamal}

\begin{itemize}
    \item Cipherteks pada EC-Elgamal adalah pasangan titik
    \begin{itemize}
        \item $P_C = ( (kB), (P_M + kP_B) )$ \quad (ket: $P_b = \text{kunci publik Bob}$)
    \end{itemize}
    
    \item $\color{red}{\text{Cipherteks pada Elgamal juga pasangan nilai:}}$
    \begin{itemize}
        \item $\color{red}{C = (a, b) = (g^k \pmod p, my_B^k \pmod p)$ \quad (ket: $y_b = \text{kunci publik Bob}$)}
    \end{itemize}
\end{itemize}

\hrulefill

\begin{itemize}
    \item Pada EC\_Elgamal, Bob mengurangkan titik kedua dari $P_C$ dengan hasil kali $b \cdot (kB)$
    \begin{itemize}
        \item $(P_M + kP_B) - [b(kB)] = P_M + k(bB) - b(kB) = P_M$
    \end{itemize}
    
    \item $\color{red}{\text{Pada El Gamal, Bob membagi nilai kedua dengan nilai pertama yang dipangkatkan}}$
    \color{red}{\text{dengan kunci privat Bob}}
    \begin{itemize}
        \item $\color{red}{my_B^k / (g^k)^b = mg^{k*b} / g^{k*b} = m}$
    \end{itemize}
\end{itemize}

\vspace{10mm}

\begin{flushright}
\small $\text{*)} \text{Sumber bahan: } \text{\color{blue}{Debdeep Mukhopadhyay, Elliptic Curve Cryptography,}}$ \\
\small $\text{\color{blue}{Dept of Computer Sc and Engg IIT Madras}}$
\end{flushright}

\end{document}
